% Options for packages loaded elsewhere
\PassOptionsToPackage{unicode}{hyperref}
\PassOptionsToPackage{hyphens}{url}
\documentclass[
  11pt,
  a4paper,
]{article}
\usepackage{xcolor}
\usepackage[margin=18mm]{geometry}
\usepackage{amsmath,amssymb}
\setcounter{secnumdepth}{-\maxdimen} % remove section numbering
\usepackage{iftex}
\ifPDFTeX
  \usepackage[T1]{fontenc}
  \usepackage[utf8]{inputenc}
  \usepackage{textcomp} % provide euro and other symbols
\else % if luatex or xetex
  \usepackage{unicode-math} % this also loads fontspec
  \defaultfontfeatures{Scale=MatchLowercase}
  \defaultfontfeatures[\rmfamily]{Ligatures=TeX,Scale=1}
\fi
\ifPDFTeX\else
  % xetex/luatex font selection
    \setmainfont[]{Noto Serif}
    \setsansfont[]{Noto Sans}
    \setmonofont[]{Noto Sans Mono}
  \setmathfont[]{Noto Sans Math}
\fi
% Use upquote if available, for straight quotes in verbatim environments
\IfFileExists{upquote.sty}{\usepackage{upquote}}{}
\IfFileExists{microtype.sty}{% use microtype if available
  \usepackage[]{microtype}
  \UseMicrotypeSet[protrusion]{basicmath} % disable protrusion for tt fonts
}{}
\makeatletter
\@ifundefined{KOMAClassName}{% if non-KOMA class
  \IfFileExists{parskip.sty}{%
    \usepackage{parskip}
  }{% else
    \setlength{\parindent}{0pt}
    \setlength{\parskip}{6pt plus 2pt minus 1pt}}
}{% if KOMA class
  \KOMAoptions{parskip=half}}
\makeatother
\ifLuaTeX
\usepackage[bidi=basic]{babel}
\else
\usepackage[bidi=default]{babel}
\fi
\babelprovide[main,import]{russian}
\ifPDFTeX
\else
\babelfont{rm}[]{Noto Serif}
\fi
% get rid of language-specific shorthands (see #6817):
\let\LanguageShortHands\languageshorthands
\def\languageshorthands#1{}
\setlength{\emergencystretch}{3em} % prevent overfull lines
\providecommand{\tightlist}{%
  \setlength{\itemsep}{0pt}\setlength{\parskip}{0pt}}
\usepackage{bookmark}
\IfFileExists{xurl.sty}{\usepackage{xurl}}{} % add URL line breaks if available
\urlstyle{same}
\hypersetup{
  pdftitle={Билеты по дискретным случайным процессам},
  pdflang={ru-RU},
  hidelinks,
  pdfcreator={LaTeX via pandoc}}

\title{Билеты по дискретным случайным процессам}
\author{}
\date{}

\begin{document}
\maketitle

\section{Билет 10. Теорема Радона-Никодима, плотности распределений и
условные математические
ожидания}\label{ux431ux438ux43bux435ux442-10.-ux442ux435ux43eux440ux435ux43cux430-ux440ux430ux434ux43eux43dux430-ux43dux438ux43aux43eux434ux438ux43cux430-ux43fux43bux43eux442ux43dux43eux441ux442ux438-ux440ux430ux441ux43fux440ux435ux434ux435ux43bux435ux43dux438ux439-ux438-ux443ux441ux43bux43eux432ux43dux44bux435-ux43cux430ux442ux435ux43cux430ux442ux438ux447ux435ux441ux43aux438ux435-ux43eux436ux438ux434ux430ux43dux438ux44f}

Источники: Ширяев 1, гл. II, §§6-7; лекции ДСП.

\subsection{1. Короткий
ответ}\label{ux43aux43eux440ux43eux442ux43aux438ux439-ux43eux442ux432ux435ux442}

Теорема Радона-Никодима говорит: если одна мера абсолютно непрерывна
относительно другой, то первая выражается через интеграл некоторой
функции по второй мере. Эта функция называется плотностью, или
производной Радона-Никодима.

Главная формула:

\[
\nu\ll\mu
\quad\Longrightarrow\quad
\nu(A)=\int_A f\,d\mu,
\qquad
f=\frac{d\nu}{d\mu}.
\]

В вероятности плотность распределения - это частный случай этой
производной:

\[
P_X\ll\lambda_n
\quad\Longrightarrow\quad
P_X(B)=\int_B p_X(x)\,dx,
\qquad
p_X=\frac{dP_X}{d\lambda_n}.
\]

Условное математическое ожидание - это среднее значение случайной
величины при заданной информации. Если \(X\in L^1\) и
\(\mathcal G\subset\mathcal F\), то

\[
Y=\mathbb E[X\mid\mathcal G]
\]

означает:

\begin{enumerate}
\def\labelenumi{\arabic{enumi}.}
\tightlist
\item
  \(Y\) является \(\mathcal G\)-измеримой;
\item
  для всех \(G\in\mathcal G\)
\end{enumerate}

\[
\int_G Y\,dP=\int_G X\,dP.
\]

Существование условного ожидания получается из теоремы Радона-Никодима:
на \(\mathcal G\) рассматривают меру со знаком

\[
Q(G)=\int_G X\,dP,
\]

она абсолютно непрерывна относительно \(P|_{\mathcal G}\), а ее
производная Радона-Никодима и есть \(\mathbb E[X\mid\mathcal G]\).

\subsection{2. Полный
ответ}\label{ux43fux43eux43bux43dux44bux439-ux43eux442ux432ux435ux442}

\subsubsection{2.1. Три уровня производной
Радона-Никодима}\label{ux442ux440ux438-ux443ux440ux43eux432ux43dux44f-ux43fux440ux43eux438ux437ux432ux43eux434ux43dux43eux439-ux440ux430ux434ux43eux43dux430-ux43dux438ux43aux43eux434ux438ux43cux430}

В этом билете соединяются три уровня одной идеи.

Первый уровень - теорема Радона-Никодима. Она отвечает на вопрос: когда
меру \(\nu\) можно записать как интеграл по другой мере \(\mu\)? Ответ:
когда \(\nu\) абсолютно непрерывна относительно \(\mu\), то есть когда
\(\mu\)-нулевые множества являются и \(\nu\)-нулевыми.

Второй уровень - плотности распределений. Если распределение случайного
вектора абсолютно непрерывно относительно меры Лебега, то у него есть
плотность. То есть плотность - это не отдельный объект, который
существует всегда, а производная Радона-Никодима распределения по мере
Лебега.

Третий уровень - условное математическое ожидание. Оно строится как
производная Радона-Никодима меры

\[
G\mapsto\int_G X\,dP
\]

относительно вероятности, но только на меньшей \(\sigma\)-алгебре
\(\mathcal G\). Поэтому условное ожидание - это не ``значение при
фиксированном исходе'', а \(\mathcal G\)-измеримая случайная величина,
которая сохраняет интегралы по всем событиям из \(\mathcal G\).

\subsubsection{2.2. Абсолютная непрерывность
мер}\label{ux430ux431ux441ux43eux43bux44eux442ux43dux430ux44f-ux43dux435ux43fux440ux435ux440ux44bux432ux43dux43eux441ux442ux44c-ux43cux435ux440}

Пусть \((E,\mathcal E)\) - измеримое пространство, а \(\mu\) и \(\nu\) -
меры на нем. Говорят, что \(\nu\) абсолютно непрерывна относительно
\(\mu\), и пишут

\[
\nu\ll\mu,
\]

если

\[
\mu(A)=0\quad\Longrightarrow\quad \nu(A)=0
\]

для всех \(A\in\mathcal E\). Интуитивно: мера \(\nu\) не видит ничего
вне того, что видит \(\mu\); у \(\nu\) нет массы на \(\mu\)-нулевых
множествах.

\subsubsection{2.3. Теорема
Радона-Никодима}\label{ux442ux435ux43eux440ux435ux43cux430-ux440ux430ux434ux43eux43dux430-ux43dux438ux43aux43eux434ux438ux43cux430}

Теорема Радона-Никодима в вероятностно нужной форме звучит так. Пусть
\(\mu\) - \(\sigma\)-конечная мера, а \(\nu\) - \(\sigma\)-конечная
неотрицательная мера на \((E,\mathcal E)\), причем

\[
\nu\ll\mu.
\]

Тогда существует \(\mathcal E\)-измеримая функция

\[
f:E\to[0,\infty]
\]

такая, что для всех \(A\in\mathcal E\)

\[
\nu(A)=\int_A f\,d\mu.
\]

Функция \(f\) называется производной Радона-Никодима меры \(\nu\)
относительно \(\mu\) и обозначается

\[
f=\frac{d\nu}{d\mu}.
\]

Эта функция единственна \(\mu\)-почти всюду. Это значит, что если \(f\)
и \(g\) дают одну и ту же меру:

\[
\nu(A)=\int_A f\,d\mu=\int_A g\,d\mu
\]

для всех \(A\), то

\[
f=g
\qquad \mu\text{-почти всюду}.
\]

Для мер со знаком формулировка аналогична: если \(\eta\) - мера со
знаком конечной вариации или более общо допустимая в стандартной
формулировке, \(\eta\ll\mu\), то существует измеримая функция \(f\),
возможно знакопеременная, такая что

\[
\eta(A)=\int_A f\,d\mu.
\]

Именно этот вариант нужен для условного ожидания, потому что

\[
Q(G)=\int_G X\,dP
\]

может быть мерой со знаком.

\subsubsection{2.4. Плотность распределения как производная
Радона-Никодима}\label{ux43fux43bux43eux442ux43dux43eux441ux442ux44c-ux440ux430ux441ux43fux440ux435ux434ux435ux43bux435ux43dux438ux44f-ux43aux430ux43a-ux43fux440ux43eux438ux437ux432ux43eux434ux43dux430ux44f-ux440ux430ux434ux43eux43dux430-ux43dux438ux43aux43eux434ux438ux43cux430}

Теперь связь с плотностями распределений. Пусть
\(X:\Omega\to\mathbb R^n\) - случайный вектор, а \(P_X\) - его
распределение:

\[
P_X(B)=P\{X\in B\},
\qquad
B\in\mathcal B(\mathbb R^n).
\]

Если

\[
P_X\ll\lambda_n,
\]

где \(\lambda_n\) - \(n\)-мерная мера Лебега, то по теореме
Радона-Никодима существует борелевская функция \(p_X\ge0\), такая что

\[
P_X(B)=\int_B p_X(x)\,\lambda_n(dx)
=
\int_B p_X(x)\,dx.
\]

Эта функция \(p_X\) называется плотностью распределения \(X\)
относительно меры Лебега:

\[
p_X=\frac{dP_X}{d\lambda_n}.
\]

Она определена не поточечно единственным образом, а только почти всюду
относительно \(\lambda_n\). Поэтому значения плотности в отдельных
точках не имеют вероятностного смысла. Вероятности получаются только
после интегрирования по множествам:

\[
P\{X\in B\}=\int_B p_X(x)\,dx.
\]

Для конечномерных распределений случайного процесса \((X_t)_{t\in T}\)
все то же самое применяется к вектору

\[
(X_{t_1},\ldots,X_{t_n}).
\]

Если его распределение абсолютно непрерывно относительно \(\lambda_n\),
то есть плотность конечномерного распределения:

\[
p_{t_1,\ldots,t_n}(x_1,\ldots,x_n)
=
\frac{dP_{(X_{t_1},\ldots,X_{t_n})}}{d\lambda_n}.
\]

Тогда

\[
P\{(X_{t_1},\ldots,X_{t_n})\in B\}
=
\int_B p_{t_1,\ldots,t_n}(x)\,dx.
\]

Если есть совместная плотность \(p_{X,Y}(x,y)\), то маргинальные
плотности находятся интегрированием:

\[
p_X(x)=\int_{\mathbb R^m}p_{X,Y}(x,y)\,dy,
\qquad
p_Y(y)=\int_{\mathbb R^n}p_{X,Y}(x,y)\,dx,
\]

что связано с Билетом 09 о Фубини.

Для условных плотностей используется та же логика. Если \((X,Y)\) имеет
совместную плотность \(p_{X,Y}(x,y)\), то для \(P_Y\)-почти всех \(y\),
для которых

\[
p_Y(y)>0,
\]

можно взять версию условной плотности \(X\) при \(Y=y\) по формуле

\[
p_{X\mid Y}(x\mid y)
=
\frac{p_{X,Y}(x,y)}{p_Y(y)}.
\]

При \(p_Y(y)=0\) значение версии условной плотности задают произвольно.

Тогда

\[
\mathbb E[g(X)\mid Y=y]
=
\int g(x)p_{X\mid Y}(x\mid y)\,dx,
\]

если интеграл определен. А условное ожидание относительно случайной
величины \(Y\) записывается как

\[
\mathbb E[g(X)\mid Y]
=
m(Y),
\]

где

\[
m(y)=\int g(x)p_{X\mid Y}(x\mid y)\,dx.
\]

Важно: формула с условной плотностью - удобный вычислительный случай.
Общий объект \(\mathbb E[X\mid\mathcal G]\) существует и без плотностей.

\subsubsection{2.5. Определение условного математического
ожидания}\label{ux43eux43fux440ux435ux434ux435ux43bux435ux43dux438ux435-ux443ux441ux43bux43eux432ux43dux43eux433ux43e-ux43cux430ux442ux435ux43cux430ux442ux438ux447ux435ux441ux43aux43eux433ux43e-ux43eux436ux438ux434ux430ux43dux438ux44f}

Перейдем к строгому определению условного математического ожидания.
Пусть

\[
(\Omega,\mathcal F,P)
\]

вероятностное пространство, \(\mathcal G\subset\mathcal F\) -
под-\(\sigma\)-алгебра, а \(X\) - интегрируемая случайная величина:

\[
\mathbb E|X|<\infty.
\]

Случайная величина \(Y\) называется условным математическим ожиданием
\(X\) относительно \(\mathcal G\), если:

\begin{enumerate}
\def\labelenumi{\arabic{enumi}.}
\tightlist
\item
  \(Y\) является \(\mathcal G\)-измеримой;
\item
  \(Y\) интегрируема;
\item
  для каждого события \(G\in\mathcal G\)
\end{enumerate}

\[
\int_G Y\,dP=\int_G X\,dP.
\]

Пишут

\[
Y=\mathbb E[X\mid\mathcal G].
\]

Смысл условий такой. \(\mathcal G\) - это доступная информация.
\(\mathcal G\)-измеримость означает, что \(Y\) зависит только от этой
информации. Интегральное равенство означает, что на любом событии,
различимом с помощью \(\mathcal G\), средняя масса \(Y\) такая же, как у
исходной величины \(X\).

\subsubsection{2.6. Существование и единственность условного
ожидания}\label{ux441ux443ux449ux435ux441ux442ux432ux43eux432ux430ux43dux438ux435-ux438-ux435ux434ux438ux43dux441ux442ux432ux435ux43dux43dux43eux441ux442ux44c-ux443ux441ux43bux43eux432ux43dux43eux433ux43e-ux43eux436ux438ux434ux430ux43dux438ux44f}

Существование доказывается через Радона-Никодима. На \(\mathcal G\)
определим функцию множеств

\[
Q(G)=\int_G X\,dP,
\qquad
G\in\mathcal G.
\]

Так как \(X\in L^1\), \(Q\) является конечной мерой со знаком на
\((\Omega,\mathcal G)\). Если

\[
P(G)=0,
\]

то

\[
Q(G)=\int_G X\,dP=0,
\]

значит

\[
Q\ll P|_{\mathcal G}.
\]

По теореме Радона-Никодима существует \(\mathcal G\)-измеримая функция
\(Y\), такая что

\[
Q(G)=\int_G Y\,dP.
\]

То есть

\[
\int_G X\,dP=\int_G Y\,dP
\]

для всех \(G\in\mathcal G\). Именно эта \(Y\) и есть
\(\mathbb E[X\mid\mathcal G]\).

Единственность условного ожидания тоже понимается почти наверное. Если
\(Y_1\) и \(Y_2\) удовлетворяют определению, то

\[
Y_1=Y_2
\qquad P\text{-почти наверное}.
\]

Поэтому условное ожидание - это класс эквивалентности случайных величин
с точностью до изменения на множестве вероятности ноль. Конкретно
выбранную функцию называют версией условного ожидания.

\subsubsection{2.7. Условная вероятность и ожидание относительно
случайной
величины}\label{ux443ux441ux43bux43eux432ux43dux430ux44f-ux432ux435ux440ux43eux44fux442ux43dux43eux441ux442ux44c-ux438-ux43eux436ux438ux434ux430ux43dux438ux435-ux43eux442ux43dux43eux441ux438ux442ux435ux43bux44cux43dux43e-ux441ux43bux443ux447ux430ux439ux43dux43eux439-ux432ux435ux43bux438ux447ux438ux43dux44b}

Условная вероятность - частный случай условного ожидания. Для события
\(A\in\mathcal F\):

\[
P(A\mid\mathcal G)=\mathbb E[\mathbf 1_A\mid\mathcal G].
\]

Это \(\mathcal G\)-измеримая случайная величина, для которой

\[
\int_G P(A\mid\mathcal G)\,dP
=
P(A\cap G),
\qquad
G\in\mathcal G.
\]

Если \(\mathcal G=\sigma(Y)\), то пишут

\[
\mathbb E[X\mid Y]
\]

вместо

\[
\mathbb E[X\mid\sigma(Y)].
\]

Так как \(\mathbb E[X\mid Y]\) является \(\sigma(Y)\)-измеримой, при
стандартных условиях она представляется как функция от \(Y\):

\[
\mathbb E[X\mid Y]=m(Y)
\]

для некоторой борелевской функции \(m\). Эту функцию часто обозначают

\[
m(y)=\mathbb E[X\mid Y=y].
\]

Но нужно помнить: \(\mathbb E[X\mid Y=y]\) как поточечная функция
определена только с точностью до \(P_Y\)-нулевых множеств. При
непрерывном \(Y\) событие \(\{Y=y\}\) обычно имеет вероятность ноль,
поэтому нельзя определять условное ожидание простой формулой
\(\mathbb E[X\mathbf 1_{\{Y=y\}}]/P\{Y=y\}\).

\subsubsection{2.8. Свойства условного
ожидания}\label{ux441ux432ux43eux439ux441ux442ux432ux430-ux443ux441ux43bux43eux432ux43dux43eux433ux43e-ux43eux436ux438ux434ux430ux43dux438ux44f}

Основные свойства условного ожидания следуют из определения и
единственности.

Линейность:

\[
\mathbb E[aX+bZ\mid\mathcal G]
=
a\mathbb E[X\mid\mathcal G]
+b\mathbb E[Z\mid\mathcal G],
\]

если все ожидания корректно определены.

Монотонность:

\[
X\le Z\ \text{п.н.}
\quad\Longrightarrow\quad
\mathbb E[X\mid\mathcal G]\le\mathbb E[Z\mid\mathcal G]\ \text{п.н.}
\]

Сохранение среднего:

\[
\mathbb E\,\mathbb E[X\mid\mathcal G]=\mathbb EX.
\]

Если \(X\) уже \(\mathcal G\)-измерима, то

\[
\mathbb E[X\mid\mathcal G]=X.
\]

Если \(X\) независима от \(\mathcal G\), то

\[
\mathbb E[X\mid\mathcal G]=\mathbb EX.
\]

Башенное свойство: если
\(\mathcal H\subset\mathcal G\subset\mathcal F\), то

\[
\mathbb E[\mathbb E[X\mid\mathcal G]\mid\mathcal H]
=
\mathbb E[X\mid\mathcal H].
\]

В частности,

\[
\mathbb E[\mathbb E[X\mid\mathcal G]]
=
\mathbb EX.
\]

Вынос известного множителя: если \(Z\) ограничена и
\(\mathcal G\)-измерима, то

\[
\mathbb E[ZX\mid\mathcal G]
=
Z\,\mathbb E[X\mid\mathcal G].
\]

Более общо формула верна при интегрируемости \(ZX\) и
\(Z\,\mathbb E[X\mid\mathcal G]\).

Это свойство часто формулируют как ``известное относительно
\(\mathcal G\) можно выносить за знак условного ожидания''.

\subsubsection{2.9. Геометрический смысл и
итоги}\label{ux433ux435ux43eux43cux435ux442ux440ux438ux447ux435ux441ux43aux438ux439-ux441ux43cux44bux441ux43b-ux438-ux438ux442ux43eux433ux438}

Есть и геометрический смысл: если \(X\in L^2\), то
\(\mathbb E[X\mid\mathcal G]\) - ортогональная проекция \(X\) на
подпространство \(\mathcal G\)-измеримых квадратично интегрируемых
случайных величин. Это связывает билет с Билетом 12.

Итак, Радон-Никодим дает общий язык плотностей: плотность меры
относительно меры, плотность распределения относительно Лебега, условная
плотность как отношение совместной и маргинальной плотностей. Условное
ожидание - тот же механизм, примененный к мере \(G\mapsto\int_G X\,dP\)
на под-\(\sigma\)-алгебре.

\subsection{3. Основные
определения}\label{ux43eux441ux43dux43eux432ux43dux44bux435-ux43eux43fux440ux435ux434ux435ux43bux435ux43dux438ux44f}

\subsubsection{Абсолютная непрерывность
мер}\label{ux430ux431ux441ux43eux43bux44eux442ux43dux430ux44f-ux43dux435ux43fux440ux435ux440ux44bux432ux43dux43eux441ux442ux44c-ux43cux435ux440-1}

Мера \(\nu\) абсолютно непрерывна относительно меры \(\mu\), если

\[
\mu(A)=0\quad\Rightarrow\quad \nu(A)=0
\]

для всех измеримых \(A\). Обозначение:

\[
\nu\ll\mu.
\]

\subsubsection{Сингулярность
мер}\label{ux441ux438ux43dux433ux443ux43bux44fux440ux43dux43eux441ux442ux44c-ux43cux435ux440}

Меры \(\nu\) и \(\mu\) сингулярны, если существует измеримое множество
\(A\), такое что

\[
\mu(A)=0,\qquad \nu(E\setminus A)=0.
\]

Иными словами, меры сосредоточены на непересекающихся с точностью до
нулевых множеств частях пространства.

\subsubsection{Производная
Радона-Никодима}\label{ux43fux440ux43eux438ux437ux432ux43eux434ux43dux430ux44f-ux440ux430ux434ux43eux43dux430-ux43dux438ux43aux43eux434ux438ux43cux430}

Если \(\nu\ll\mu\) и

\[
\nu(A)=\int_A f\,d\mu
\]

для всех \(A\), то

\[
f=\frac{d\nu}{d\mu}
\]

называется производной Радона-Никодима, или плотностью \(\nu\)
относительно \(\mu\).

\subsubsection{Плотность
распределения}\label{ux43fux43bux43eux442ux43dux43eux441ux442ux44c-ux440ux430ux441ux43fux440ux435ux434ux435ux43bux435ux43dux438ux44f}

Если \(P_X\ll\lambda_n\), то

\[
p_X=\frac{dP_X}{d\lambda_n}
\]

называется плотностью распределения \(X\) относительно меры Лебега.

\subsubsection{Плотность конечномерного
распределения}\label{ux43fux43bux43eux442ux43dux43eux441ux442ux44c-ux43aux43eux43dux435ux447ux43dux43eux43cux435ux440ux43dux43eux433ux43e-ux440ux430ux441ux43fux440ux435ux434ux435ux43bux435ux43dux438ux44f}

Для процесса \((X_t)\) плотность вектора \((X_{t_1},\ldots,X_{t_n})\):

\[
p_{t_1,\ldots,t_n}
=
\frac{dP_{(X_{t_1},\ldots,X_{t_n})}}{d\lambda_n},
\]

если соответствующее распределение абсолютно непрерывно относительно
\(\lambda_n\).

\subsubsection{\texorpdfstring{Под-\(\sigma\)-алгебра}{Под-\textbackslash sigma-алгебра}}\label{ux43fux43eux434-sigma-ux430ux43bux433ux435ux431ux440ux430}

\(\mathcal G\subset\mathcal F\) - \(\sigma\)-алгебра, содержащая
доступную информацию. Чем больше \(\mathcal G\), тем больше событий мы
различаем.

\subsubsection{\texorpdfstring{Условное математическое ожидание
относительно
\(\mathcal G\)}{Условное математическое ожидание относительно \textbackslash mathcal G}}\label{ux443ux441ux43bux43eux432ux43dux43eux435-ux43cux430ux442ux435ux43cux430ux442ux438ux447ux435ux441ux43aux43eux435-ux43eux436ux438ux434ux430ux43dux438ux435-ux43eux442ux43dux43eux441ux438ux442ux435ux43bux44cux43dux43e-mathcal-g}

Для \(X\in L^1\):

\[
Y=\mathbb E[X\mid\mathcal G]
\]

если \(Y\) \(\mathcal G\)-измерима, интегрируема и

\[
\int_GY\,dP=\int_GX\,dP
\]

для всех \(G\in\mathcal G\).

\subsubsection{Версия условного
ожидания}\label{ux432ux435ux440ux441ux438ux44f-ux443ux441ux43bux43eux432ux43dux43eux433ux43e-ux43eux436ux438ux434ux430ux43dux438ux44f}

Конкретная \(\mathcal G\)-измеримая функция \(Y\), удовлетворяющая
определению. Две версии совпадают \(P\)-почти наверное.

\subsubsection{\texorpdfstring{Условная вероятность относительно
\(\mathcal G\)}{Условная вероятность относительно \textbackslash mathcal G}}\label{ux443ux441ux43bux43eux432ux43dux430ux44f-ux432ux435ux440ux43eux44fux442ux43dux43eux441ux442ux44c-ux43eux442ux43dux43eux441ux438ux442ux435ux43bux44cux43dux43e-mathcal-g}

\[
P(A\mid\mathcal G)=\mathbb E[\mathbf 1_A\mid\mathcal G].
\]

\subsubsection{Условное ожидание относительно случайной
величины}\label{ux443ux441ux43bux43eux432ux43dux43eux435-ux43eux436ux438ux434ux430ux43dux438ux435-ux43eux442ux43dux43eux441ux438ux442ux435ux43bux44cux43dux43e-ux441ux43bux443ux447ux430ux439ux43dux43eux439-ux432ux435ux43bux438ux447ux438ux43dux44b}

\[
\mathbb E[X\mid Y]
=
\mathbb E[X\mid\sigma(Y)].
\]

Обычно

\[
\mathbb E[X\mid Y]=m(Y)
\]

для некоторой борелевской функции \(m\).

\subsubsection{Условная
плотность}\label{ux443ux441ux43bux43eux432ux43dux430ux44f-ux43fux43bux43eux442ux43dux43eux441ux442ux44c}

Если \((X,Y)\) имеет совместную плотность \(p_{X,Y}\) и \(Y\) имеет
плотность \(p_Y\), то при \(p_Y(y)>0\)

\[
p_{X\mid Y}(x\mid y)
=
\frac{p_{X,Y}(x,y)}{p_Y(y)}.
\]

\subsection{4. Основные теоремы и
свойства}\label{ux43eux441ux43dux43eux432ux43dux44bux435-ux442ux435ux43eux440ux435ux43cux44b-ux438-ux441ux432ux43eux439ux441ux442ux432ux430}

\subsubsection{Теорема 1. Радона-Никодима для неотрицательных
мер}\label{ux442ux435ux43eux440ux435ux43cux430-1.-ux440ux430ux434ux43eux43dux430-ux43dux438ux43aux43eux434ux438ux43cux430-ux434ux43bux44f-ux43dux435ux43eux442ux440ux438ux446ux430ux442ux435ux43bux44cux43dux44bux445-ux43cux435ux440}

Пусть \(\mu\) - \(\sigma\)-конечная мера, \(\nu\) - \(\sigma\)-конечная
неотрицательная мера и

\[
\nu\ll\mu.
\]

Тогда существует измеримая функция \(f\ge0\), такая что

\[
\nu(A)=\int_A f\,d\mu
\]

для всех \(A\). Функция \(f\) единственна \(\mu\)-почти всюду и
обозначается

\[
f=\frac{d\nu}{d\mu}.
\]

\subsubsection{Теорема 2. Радона-Никодима для меры со
знаком}\label{ux442ux435ux43eux440ux435ux43cux430-2.-ux440ux430ux434ux43eux43dux430-ux43dux438ux43aux43eux434ux438ux43cux430-ux434ux43bux44f-ux43cux435ux440ux44b-ux441ux43e-ux437ux43dux430ux43aux43eux43c}

Если \(\eta\) - конечная мера со знаком, \(\eta\ll\mu\), а \(\mu\)
\(\sigma\)-конечна, то существует измеримая функция \(f\), интегрируемая
по \(\mu\), такая что

\[
\eta(A)=\int_A f\,d\mu.
\]

Этот вариант нужен для построения \(\mathbb E[X\mid\mathcal G]\) при
знакопеременной \(X\).

\subsubsection{Теорема 3. Плотность распределения как производная
Радона-Никодима}\label{ux442ux435ux43eux440ux435ux43cux430-3.-ux43fux43bux43eux442ux43dux43eux441ux442ux44c-ux440ux430ux441ux43fux440ux435ux434ux435ux43bux435ux43dux438ux44f-ux43aux430ux43a-ux43fux440ux43eux438ux437ux432ux43eux434ux43dux430ux44f-ux440ux430ux434ux43eux43dux430-ux43dux438ux43aux43eux434ux438ux43cux430}

Для случайного вектора \(X\) в \(\mathbb R^n\) плотность существует
тогда и только тогда, когда

\[
P_X\ll\lambda_n.
\]

Тогда

\[
p_X=\frac{dP_X}{d\lambda_n},
\qquad
P\{X\in B\}=\int_B p_X(x)\,dx.
\]

\subsubsection{Свойство 4. Плотность определена почти
всюду}\label{ux441ux432ux43eux439ux441ux442ux432ux43e-4.-ux43fux43bux43eux442ux43dux43eux441ux442ux44c-ux43eux43fux440ux435ux434ux435ux43bux435ux43dux430-ux43fux43eux447ux442ux438-ux432ux441ux44eux434ux443}

Если \(p\) и \(\tilde p\) - две плотности одного распределения, то

\[
p=\tilde p
\qquad \lambda_n\text{-почти всюду}.
\]

Значения плотности в отдельных точках можно менять без изменения
распределения.

\subsubsection{Теорема 5. Существование условного математического
ожидания}\label{ux442ux435ux43eux440ux435ux43cux430-5.-ux441ux443ux449ux435ux441ux442ux432ux43eux432ux430ux43dux438ux435-ux443ux441ux43bux43eux432ux43dux43eux433ux43e-ux43cux430ux442ux435ux43cux430ux442ux438ux447ux435ux441ux43aux43eux433ux43e-ux43eux436ux438ux434ux430ux43dux438ux44f}

Если \(X\in L^1\) и \(\mathcal G\subset\mathcal F\), то существует

\[
\mathbb E[X\mid\mathcal G],
\]

причем оно единственно \(P\)-почти наверное.

\subsubsection{Доказательство}\label{ux434ux43eux43aux430ux437ux430ux442ux435ux43bux44cux441ux442ux432ux43e}

На \((\Omega,\mathcal G)\) задаем меру со знаком

\[
Q(G)=\int_GX\,dP.
\]

Она абсолютно непрерывна относительно \(P|_{\mathcal G}\). По
Радону-Никодиму существует \(\mathcal G\)-измеримая \(Y\), такая что

\[
Q(G)=\int_GY\,dP.
\]

Это и есть определение условного ожидания.

\subsubsection{Свойство 6.
Линейность}\label{ux441ux432ux43eux439ux441ux442ux432ux43e-6.-ux43bux438ux43dux435ux439ux43dux43eux441ux442ux44c}

\[
\mathbb E[aX+bZ\mid\mathcal G]
=
a\mathbb E[X\mid\mathcal G]
+b\mathbb E[Z\mid\mathcal G].
\]

\subsubsection{Свойство 7.
Монотонность}\label{ux441ux432ux43eux439ux441ux442ux432ux43e-7.-ux43cux43eux43dux43eux442ux43eux43dux43dux43eux441ux442ux44c}

\[
X\le Z\ \text{п.н.}
\quad\Rightarrow\quad
\mathbb E[X\mid\mathcal G]\le \mathbb E[Z\mid\mathcal G]\ \text{п.н.}
\]

\subsubsection{Свойство 8. Сохранение
среднего}\label{ux441ux432ux43eux439ux441ux442ux432ux43e-8.-ux441ux43eux445ux440ux430ux43dux435ux43dux438ux435-ux441ux440ux435ux434ux43dux435ux433ux43e}

\[
\mathbb E\,\mathbb E[X\mid\mathcal G]=\mathbb EX.
\]

\subsubsection{Свойство 9. Если величина уже
известна}\label{ux441ux432ux43eux439ux441ux442ux432ux43e-9.-ux435ux441ux43bux438-ux432ux435ux43bux438ux447ux438ux43dux430-ux443ux436ux435-ux438ux437ux432ux435ux441ux442ux43dux430}

Если \(X\) \(\mathcal G\)-измерима, то

\[
\mathbb E[X\mid\mathcal G]=X.
\]

\subsubsection{Свойство 10.
Независимость}\label{ux441ux432ux43eux439ux441ux442ux432ux43e-10.-ux43dux435ux437ux430ux432ux438ux441ux438ux43cux43eux441ux442ux44c}

Если \(X\) независима от \(\mathcal G\), то

\[
\mathbb E[X\mid\mathcal G]=\mathbb EX.
\]

\subsubsection{Свойство 11. Башенное
свойство}\label{ux441ux432ux43eux439ux441ux442ux432ux43e-11.-ux431ux430ux448ux435ux43dux43dux43eux435-ux441ux432ux43eux439ux441ux442ux432ux43e}

Если \(\mathcal H\subset\mathcal G\), то

\[
\mathbb E[\mathbb E[X\mid\mathcal G]\mid\mathcal H]
=
\mathbb E[X\mid\mathcal H].
\]

\subsubsection{Свойство 12. Вынос известного
множителя}\label{ux441ux432ux43eux439ux441ux442ux432ux43e-12.-ux432ux44bux43dux43eux441-ux438ux437ux432ux435ux441ux442ux43dux43eux433ux43e-ux43cux43dux43eux436ux438ux442ux435ux43bux44f}

Если \(Z\) \(\mathcal G\)-измерима и \(ZX\in L^1\), то

\[
\mathbb E[ZX\mid\mathcal G]
=
Z\mathbb E[X\mid\mathcal G].
\]

\subsubsection{Свойство 13. Условная
вероятность}\label{ux441ux432ux43eux439ux441ux442ux432ux43e-13.-ux443ux441ux43bux43eux432ux43dux430ux44f-ux432ux435ux440ux43eux44fux442ux43dux43eux441ux442ux44c}

Для \(A\in\mathcal F\):

\[
P(A\mid\mathcal G)=\mathbb E[\mathbf 1_A\mid\mathcal G],
\]

и для всех \(G\in\mathcal G\)

\[
\int_G P(A\mid\mathcal G)\,dP=P(A\cap G).
\]

\subsubsection{Свойство 14. Условное ожидание при совместной
плотности}\label{ux441ux432ux43eux439ux441ux442ux432ux43e-14.-ux443ux441ux43bux43eux432ux43dux43eux435-ux43eux436ux438ux434ux430ux43dux438ux435-ux43fux440ux438-ux441ux43eux432ux43cux435ux441ux442ux43dux43eux439-ux43fux43bux43eux442ux43dux43eux441ux442ux438}

Если \((X,Y)\) имеет совместную плотность \(p_{X,Y}\), то

\[
\mathbb E[g(X)\mid Y=y]
=
\int g(x)\frac{p_{X,Y}(x,y)}{p_Y(y)}\,dx
\]

при \(p_Y(y)>0\). Тогда

\[
\mathbb E[g(X)\mid Y]=m(Y),
\]

где

\[
m(y)=\int g(x)p_{X\mid Y}(x\mid y)\,dx.
\]

\subsection{5. Примеры}\label{ux43fux440ux438ux43cux435ux440ux44b}

\subsubsection{Пример 1. Обычная плотность как
Радон-Никодим}\label{ux43fux440ux438ux43cux435ux440-1.-ux43eux431ux44bux447ux43dux430ux44f-ux43fux43bux43eux442ux43dux43eux441ux442ux44c-ux43aux430ux43a-ux440ux430ux434ux43eux43d-ux43dux438ux43aux43eux434ux438ux43c}

Пусть \(X\) равномерна на \([0,1]\). Тогда

\[
P_X(B)=\lambda(B\cap[0,1]).
\]

Значит

\[
P_X\ll\lambda,
\qquad
\frac{dP_X}{d\lambda}(x)=\mathbf 1_{[0,1]}(x)
\]

почти всюду.

\subsubsection{Пример 2. Нет плотности относительно
Лебега}\label{ux43fux440ux438ux43cux435ux440-2.-ux43dux435ux442-ux43fux43bux43eux442ux43dux43eux441ux442ux438-ux43eux442ux43dux43eux441ux438ux442ux435ux43bux44cux43dux43e-ux43bux435ux431ux435ux433ux430}

Если \(X\equiv a\), то

\[
P_X=\delta_a.
\]

Мера \(\delta_a\) не абсолютно непрерывна относительно меры Лебега,
потому что

\[
\lambda(\{a\})=0,
\qquad
\delta_a(\{a\})=1.
\]

Значит плотности относительно Лебега нет. Это важный контрпример:
распределение есть всегда, плотность - не всегда.

\subsubsection{Пример 3. Условное ожидание по конечному
разбиению}\label{ux43fux440ux438ux43cux435ux440-3.-ux443ux441ux43bux43eux432ux43dux43eux435-ux43eux436ux438ux434ux430ux43dux438ux435-ux43fux43e-ux43aux43eux43dux435ux447ux43dux43eux43cux443-ux440ux430ux437ux431ux438ux435ux43dux438ux44e}

Пусть

\[
\mathcal G=\sigma(D_1,\ldots,D_n),
\]

где \(D_i\) попарно не пересекаются, \(\bigcup_iD_i=\Omega\),
\(P(D_i)>0\). Тогда

\[
\mathbb E[X\mid\mathcal G]
=
\sum_{i=1}^n
\frac{\int_{D_i}X\,dP}{P(D_i)}\mathbf 1_{D_i}.
\]

То есть на каждом атоме разбиения условное ожидание равно среднему \(X\)
по этому атому.

\subsubsection{Пример 4. Условное ожидание при
независимости}\label{ux43fux440ux438ux43cux435ux440-4.-ux443ux441ux43bux43eux432ux43dux43eux435-ux43eux436ux438ux434ux430ux43dux438ux435-ux43fux440ux438-ux43dux435ux437ux430ux432ux438ux441ux438ux43cux43eux441ux442ux438}

Если \(X\) независима от \(\mathcal G\), то

\[
\mathbb E[X\mid\mathcal G]=\mathbb EX.
\]

Например, если \(X\) и \(Y\) независимы, то

\[
\mathbb E[X\mid Y]=\mathbb EX.
\]

\subsubsection{Пример 5. Если величина уже
известна}\label{ux43fux440ux438ux43cux435ux440-5.-ux435ux441ux43bux438-ux432ux435ux43bux438ux447ux438ux43dux430-ux443ux436ux435-ux438ux437ux432ux435ux441ux442ux43dux430}

Если \(X\) является функцией от \(Y\), например

\[
X=Y^2,
\]

то \(X\) \(\sigma(Y)\)-измерима, поэтому

\[
\mathbb E[X\mid Y]=X=Y^2.
\]

\subsubsection{Пример 6. Условное ожидание через совместную
плотность}\label{ux43fux440ux438ux43cux435ux440-6.-ux443ux441ux43bux43eux432ux43dux43eux435-ux43eux436ux438ux434ux430ux43dux438ux435-ux447ux435ux440ux435ux437-ux441ux43eux432ux43cux435ux441ux442ux43dux443ux44e-ux43fux43bux43eux442ux43dux43eux441ux442ux44c}

Пусть \((X,Y)\) имеет совместную плотность \(p_{X,Y}\). Тогда

\[
p_Y(y)=\int p_{X,Y}(x,y)\,dx.
\]

Если \(p_Y(y)>0\), то

\[
\mathbb E[X\mid Y=y]
=
\int x\frac{p_{X,Y}(x,y)}{p_Y(y)}\,dx.
\]

В частности,

\[
\mathbb E[X\mid Y]=m(Y),
\]

где

\[
m(y)=\int x\,p_{X\mid Y}(x\mid y)\,dx.
\]

\subsubsection{Пример 7. Условное ожидание как
прогноз}\label{ux43fux440ux438ux43cux435ux440-7.-ux443ux441ux43bux43eux432ux43dux43eux435-ux43eux436ux438ux434ux430ux43dux438ux435-ux43aux430ux43a-ux43fux440ux43eux433ux43dux43eux437}

Пусть

\[
X=Y+\varepsilon,
\]

где \(\varepsilon\) независима от \(Y\) и \(\mathbb E\varepsilon=0\).
Тогда

\[
\mathbb E[X\mid Y]
=
\mathbb E[Y+\varepsilon\mid Y]
=
Y+\mathbb E[\varepsilon\mid Y]
=
Y.
\]

Здесь \(Y\) - известная часть, а независимый шум усредняется до нуля.

\subsection{6. Схемы доказательств /
обоснований}\label{ux441ux445ux435ux43cux44b-ux434ux43eux43aux430ux437ux430ux442ux435ux43bux44cux441ux442ux432-ux43eux431ux43eux441ux43dux43eux432ux430ux43dux438ux439}

\textbf{Идея теоремы Радона-Никодима.}

Полное доказательство - теорема теории меры. Интуитивная идея такая:
если \(\nu\ll\mu\), то \(\nu\) не имеет массы вне \(\mu\)-массы, значит
ее можно описать локальной ``плотностью'' относительно \(\mu\). В
простом случае, когда пространство разбито на конечные атомы \(A_i\) с
\(\mu(A_i)>0\), плотность равна

\[
f=\sum_i \frac{\nu(A_i)}{\mu(A_i)}\mathbf 1_{A_i}.
\]

Общая теорема строит аналог этой локальной плотности для произвольной
\(\sigma\)-алгебры.

\textbf{Доказательство существования условного ожидания.}

\begin{enumerate}
\def\labelenumi{\arabic{enumi}.}
\tightlist
\item
  Берем \(X\in L^1\) и под-\(\sigma\)-алгебру \(\mathcal G\).
\item
  На \(\mathcal G\) задаем
\end{enumerate}

\[
Q(G)=\int_GX\,dP.
\]

\begin{enumerate}
\def\labelenumi{\arabic{enumi}.}
\setcounter{enumi}{2}
\tightlist
\item
  \(Q\) - конечная мера со знаком.
\item
  Если \(P(G)=0\), то \(Q(G)=0\), значит \(Q\ll P|_{\mathcal G}\).
\item
  По Радону-Никодиму существует \(Y=dQ/dP|_{\mathcal G}\).
\item
  Получаем
\end{enumerate}

\[
\int_GY\,dP=Q(G)=\int_GX\,dP.
\]

Значит \(Y=\mathbb E[X\mid\mathcal G]\).

\textbf{Доказательство единственности условного ожидания.}

Пусть \(Y_1\) и \(Y_2\) удовлетворяют определению. Тогда для всех
\(G\in\mathcal G\)

\[
\int_G(Y_1-Y_2)\,dP=0.
\]

Положим

\[
G_+=\{Y_1>Y_2\}.
\]

Так как \(Y_1,Y_2\) \(\mathcal G\)-измеримы, то \(G_+\in\mathcal G\).
Тогда

\[
\int_{G_+}(Y_1-Y_2)\,dP=0.
\]

Но подынтегральная функция положительна на \(G_+\), значит \(P(G_+)=0\).
Аналогично для \(\{Y_2>Y_1\}\). Поэтому \(Y_1=Y_2\) почти наверное.

\textbf{Доказательство башенного свойства.}

Пусть \(\mathcal H\subset\mathcal G\). Обозначим

\[
Z=\mathbb E[\mathbb E[X\mid\mathcal G]\mid\mathcal H].
\]

Для любого \(H\in\mathcal H\) имеем \(H\in\mathcal G\), поэтому

\[
\int_H Z\,dP
=
\int_H \mathbb E[X\mid\mathcal G]\,dP
=
\int_H X\,dP.
\]

Кроме того, \(Z\) \(\mathcal H\)-измерима. Значит по определению

\[
Z=\mathbb E[X\mid\mathcal H].
\]

\textbf{Доказательство свойства независимости.}

Если \(X\) независима от \(\mathcal G\), то постоянная \(\mathbb EX\)
является \(\mathcal G\)-измеримой. Для любого \(G\in\mathcal G\):

\[
\int_G X\,dP
=
\mathbb E[X\mathbf 1_G]
=
\mathbb EX\,P(G)
=
\int_G \mathbb EX\,dP.
\]

По единственности

\[
\mathbb E[X\mid\mathcal G]=\mathbb EX.
\]

\textbf{Доказательство формулы для условной плотности.}

Для борелевского \(A\):

\[
P\{X\in A,\ Y\in dy\}
\]

имеет плотность по \(y\), равную

\[
\int_A p_{X,Y}(x,y)\,dx.
\]

Делим на маргинальную плотность

\[
p_Y(y)=\int p_{X,Y}(x,y)\,dx
\]

и получаем условную плотность

\[
p_{X\mid Y}(x\mid y)=\frac{p_{X,Y}(x,y)}{p_Y(y)}.
\]

\subsection{7. Что написать на
доске}\label{ux447ux442ux43e-ux43dux430ux43fux438ux441ux430ux442ux44c-ux43dux430-ux434ux43eux441ux43aux435}

\begin{enumerate}
\def\labelenumi{\arabic{enumi}.}
\tightlist
\item
  Абсолютная непрерывность:
\end{enumerate}

\[
\nu\ll\mu
\quad\Longleftrightarrow\quad
\mu(A)=0\Rightarrow\nu(A)=0.
\]

\begin{enumerate}
\def\labelenumi{\arabic{enumi}.}
\setcounter{enumi}{1}
\tightlist
\item
  Радон-Никодим:
\end{enumerate}

\[
\nu(A)=\int_A \frac{d\nu}{d\mu}\,d\mu.
\]

\begin{enumerate}
\def\labelenumi{\arabic{enumi}.}
\setcounter{enumi}{2}
\tightlist
\item
  Плотность распределения:
\end{enumerate}

\[
P_X\ll\lambda_n,
\qquad
p_X=\frac{dP_X}{d\lambda_n},
\qquad
P_X(B)=\int_Bp_X(x)\,dx.
\]

\begin{enumerate}
\def\labelenumi{\arabic{enumi}.}
\setcounter{enumi}{3}
\tightlist
\item
  Условное ожидание:
\end{enumerate}

\[
Y=\mathbb E[X\mid\mathcal G]
\iff
\begin{cases}
Y\ \mathcal G\text{-измерима},\\
\int_GY\,dP=\int_GX\,dP,\quad G\in\mathcal G.
\end{cases}
\]

\begin{enumerate}
\def\labelenumi{\arabic{enumi}.}
\setcounter{enumi}{4}
\tightlist
\item
  Построение через Радона-Никодима:
\end{enumerate}

\[
Q(G)=\int_GX\,dP,
\qquad
\mathbb E[X\mid\mathcal G]=\frac{dQ}{dP|_{\mathcal G}}.
\]

\begin{enumerate}
\def\labelenumi{\arabic{enumi}.}
\setcounter{enumi}{5}
\tightlist
\item
  Главные свойства:
\end{enumerate}

\[
\mathbb E\,\mathbb E[X\mid\mathcal G]=\mathbb EX,
\]

\[
\mathcal H\subset\mathcal G
\Rightarrow
\mathbb E[\mathbb E[X\mid\mathcal G]\mid\mathcal H]
=
\mathbb E[X\mid\mathcal H],
\]

\[
X\perp\mathcal G
\Rightarrow
\mathbb E[X\mid\mathcal G]=\mathbb EX.
\]

\begin{enumerate}
\def\labelenumi{\arabic{enumi}.}
\setcounter{enumi}{6}
\tightlist
\item
  Условная плотность:
\end{enumerate}

\[
p_{X\mid Y}(x\mid y)=\frac{p_{X,Y}(x,y)}{p_Y(y)}.
\]

\subsection{8. Типичные дополнительные
вопросы}\label{ux442ux438ux43fux438ux447ux43dux44bux435-ux434ux43eux43fux43eux43bux43dux438ux442ux435ux43bux44cux43dux44bux435-ux432ux43eux43fux440ux43eux441ux44b}

\textbf{Вопрос. Что означает \(\nu\ll\mu\)?}

Это значит, что всякое \(\mu\)-нулевое множество является
\(\nu\)-нулевым:

\[
\mu(A)=0\Rightarrow\nu(A)=0.
\]

\textbf{Вопрос. Почему плотность распределения - это производная
Радона-Никодима?}

Потому что плотность \(p_X\) задает распределение через интеграл по мере
Лебега:

\[
P_X(B)=\int_Bp_X\,d\lambda.
\]

Это ровно представление Радона-Никодима для \(P_X\ll\lambda\).

\textbf{Вопрос. Любое ли распределение имеет плотность?}

Нет. Распределение Дирака \(\delta_a\) не имеет плотности относительно
меры Лебега, потому что оно дает массу точке, а мера Лебега этой точки
равна нулю.

\textbf{Вопрос. Плотность единственна?}

Только почти всюду относительно меры, по которой берется плотность.
Значения на нулевых множествах можно менять.

\textbf{Вопрос. Что такое условное ожидание строго?}

Это \(\mathcal G\)-измеримая случайная величина \(Y\), такая что

\[
\int_GY\,dP=\int_GX\,dP
\]

для всех \(G\in\mathcal G\).

\textbf{Вопрос. Почему условное ожидание существует?}

Потому что мера со знаком

\[
Q(G)=\int_GX\,dP
\]

абсолютно непрерывна относительно \(P\) на \(\mathcal G\), и можно
применить теорему Радона-Никодима.

\textbf{Вопрос. Почему условное ожидание определено только почти
наверное?}

Потому что производная Радона-Никодима единственна только с точностью до
изменения на множестве меры ноль.

\textbf{Вопрос. Что будет, если \(\mathcal G\) тривиальна?}

Если

\[
\mathcal G=\{\varnothing,\Omega\},
\]

то

\[
\mathbb E[X\mid\mathcal G]=\mathbb EX.
\]

\textbf{Вопрос. Что будет, если \(X\) \(\mathcal G\)-измерима?}

Тогда

\[
\mathbb E[X\mid\mathcal G]=X.
\]

\textbf{Вопрос. Что будет, если \(X\) независима от \(\mathcal G\)?}

Тогда

\[
\mathbb E[X\mid\mathcal G]=\mathbb EX.
\]

\textbf{Вопрос. Что такое башенное свойство?}

Если \(\mathcal H\subset\mathcal G\), то

\[
\mathbb E[\mathbb E[X\mid\mathcal G]\mid\mathcal H]
=
\mathbb E[X\mid\mathcal H].
\]

\textbf{Вопрос. Чем условная вероятность отличается от условного
ожидания?}

Условная вероятность - частный случай:

\[
P(A\mid\mathcal G)=\mathbb E[\mathbf 1_A\mid\mathcal G].
\]

\textbf{Вопрос. Почему нельзя всегда писать
\(P(A\mid Y=y)=P(A\cap\{Y=y\})/P(Y=y)\)?}

Если \(Y\) непрерывна, то обычно

\[
P(Y=y)=0.
\]

Нужно использовать условное ожидание относительно \(\sigma(Y)\) или
регулярные условные распределения.

\textbf{Вопрос. Как найти \(\mathbb E[X\mid Y]\) при совместной
плотности?}

Сначала найти

\[
p_Y(y)=\int p_{X,Y}(x,y)\,dx,
\]

затем

\[
p_{X\mid Y}(x\mid y)=\frac{p_{X,Y}(x,y)}{p_Y(y)},
\]

и после этого

\[
\mathbb E[X\mid Y=y]=\int x p_{X\mid Y}(x\mid y)\,dx.
\]

\subsection{9. Типичные
ошибки}\label{ux442ux438ux43fux438ux447ux43dux44bux435-ux43eux448ux438ux431ux43aux438}

\begin{enumerate}
\def\labelenumi{\arabic{enumi}.}
\tightlist
\item
  Говорить о плотности распределения, не проверив абсолютную
  непрерывность относительно меры Лебега.
\item
  Считать значение плотности \(p(x)\) вероятностью точки.
\item
  Забывать, что производная Радона-Никодима единственна только почти
  всюду.
\item
  Путать абсолютную непрерывность меры \(\nu\ll\mu\) с абсолютной
  непрерывностью функции распределения.
\item
  Путать условное ожидание и условную вероятность: условная вероятность
  - это условное ожидание индикатора.
\item
  Определять \(\mathbb E[X\mid Y=y]\) через деление на \(P(Y=y)\) в
  непрерывном случае.
\item
  Забывать условие \(X\in L^1\) для обычного условного ожидания.
\item
  Думать, что \(\mathbb E[X\mid\mathcal G]\) - число. Обычно это
  случайная величина.
\item
  Забывать \(\mathcal G\)-измеримость в определении условного ожидания.
\item
  Применять формулу условной плотности при \(p_Y(y)=0\).
\item
  Путать равенство почти наверное с поточечным равенством.
\item
  Забывать, что при независимости условное ожидание становится
  константой \(\mathbb EX\).
\end{enumerate}

\subsection{10. Связи с другими
билетами}\label{ux441ux432ux44fux437ux438-ux441-ux434ux440ux443ux433ux438ux43cux438-ux431ux438ux43bux435ux442ux430ux43cux438}

\textbf{Билет 06.} Плотности распределений в общем виде определяются как
производные Радона-Никодима.

\textbf{Билет 07.} Все определения используют интеграл Лебега и свойства
интегрируемости.

\textbf{Билет 09.} Совместные и маргинальные плотности, условные
плотности и независимость используют Фубини.

\textbf{Билет 11.} Замена переменной позволяет записывать

\[
\mathbb E g(X)=\int g\,dP_X,
\]

а при плотности

\[
\mathbb E g(X)=\int g(x)p_X(x)\,dx.
\]

\textbf{Билет 12.} В \(L^2\) условное ожидание является ортогональным
проектором на подпространство \(\mathcal G\)-измеримых величин.

\textbf{Билет 14.} Условное ожидание связано с регрессией и оптимальным
среднеквадратичным прогнозом.

\textbf{Билет 22.} Марковские переходные функции можно понимать как
условные распределения будущего при известном текущем состоянии.

\textbf{Билет 23.} Марковское свойство формулируется через условные
вероятности и условные распределения.

\subsection{11. Минимум на
удовлетворительно}\label{ux43cux438ux43dux438ux43cux443ux43c-ux43dux430-ux443ux434ux43eux432ux43bux435ux442ux432ux43eux440ux438ux442ux435ux43bux44cux43dux43e}

Нужно уметь сказать:

\begin{enumerate}
\def\labelenumi{\arabic{enumi}.}
\tightlist
\item
  Абсолютная непрерывность:
\end{enumerate}

\[
\nu\ll\mu
\quad\Longleftrightarrow\quad
\mu(A)=0\Rightarrow\nu(A)=0.
\]

\begin{enumerate}
\def\labelenumi{\arabic{enumi}.}
\setcounter{enumi}{1}
\tightlist
\item
  Радон-Никодим:
\end{enumerate}

\[
\nu(A)=\int_A f\,d\mu,
\qquad
f=\frac{d\nu}{d\mu}.
\]

\begin{enumerate}
\def\labelenumi{\arabic{enumi}.}
\setcounter{enumi}{2}
\tightlist
\item
  Плотность распределения:
\end{enumerate}

\[
P_X(B)=\int_Bp_X(x)\,dx.
\]

\begin{enumerate}
\def\labelenumi{\arabic{enumi}.}
\setcounter{enumi}{3}
\tightlist
\item
  Условное ожидание:
\end{enumerate}

\[
Y=\mathbb E[X\mid\mathcal G]
\]

если \(Y\) \(\mathcal G\)-измерима и

\[
\int_GY\,dP=\int_GX\,dP
\]

для всех \(G\in\mathcal G\).

\begin{enumerate}
\def\labelenumi{\arabic{enumi}.}
\setcounter{enumi}{4}
\tightlist
\item
  Условная вероятность:
\end{enumerate}

\[
P(A\mid\mathcal G)=\mathbb E[\mathbf 1_A\mid\mathcal G].
\]

\subsection{12. Минимум на
хорошо}\label{ux43cux438ux43dux438ux43cux443ux43c-ux43dux430-ux445ux43eux440ux43eux448ux43e}

Помимо минимума, нужно:

\begin{enumerate}
\def\labelenumi{\arabic{enumi}.}
\tightlist
\item
  объяснить, что плотность существует не всегда;
\item
  сказать, что плотность и условное ожидание определены почти всюду;
\item
  вывести условное ожидание через Радона-Никодима:
\end{enumerate}

\[
Q(G)=\int_GX\,dP,
\qquad
\mathbb E[X\mid\mathcal G]=\frac{dQ}{dP|_{\mathcal G}};
\]

\begin{enumerate}
\def\labelenumi{\arabic{enumi}.}
\setcounter{enumi}{3}
\tightlist
\item
  назвать свойства:
\end{enumerate}

\[
\mathbb E\,\mathbb E[X\mid\mathcal G]=\mathbb EX,
\]

\[
\mathbb E[X\mid\mathcal G]=X
\quad\text{если }X\text{ }\mathcal G\text{-измерима},
\]

\[
\mathbb E[X\mid\mathcal G]=\mathbb EX
\quad\text{если }X\perp\mathcal G;
\]

\begin{enumerate}
\def\labelenumi{\arabic{enumi}.}
\setcounter{enumi}{4}
\tightlist
\item
  привести пример условного ожидания по разбиению.
\end{enumerate}

\subsection{13. Что добавить для
отлично}\label{ux447ux442ux43e-ux434ux43eux431ux430ux432ux438ux442ux44c-ux434ux43bux44f-ux43eux442ux43bux438ux447ux43dux43e}

Для сильного ответа стоит добавить:

\begin{enumerate}
\def\labelenumi{\arabic{enumi}.}
\tightlist
\item
  формулировку Радона-Никодима для мер со знаком;
\item
  доказательство существования условного ожидания через меру
  \(Q(G)=\int_GX\,dP\);
\item
  доказательство единственности почти наверное;
\item
  башенное свойство;
\item
  вынос \(\mathcal G\)-измеримого множителя;
\item
  условную плотность
\end{enumerate}

\[
p_{X\mid Y}(x\mid y)=\frac{p_{X,Y}(x,y)}{p_Y(y)};
\]

\begin{enumerate}
\def\labelenumi{\arabic{enumi}.}
\setcounter{enumi}{6}
\tightlist
\item
  предупреждение про непрерывное условие \(Y=y\), где \(P(Y=y)=0\).
\end{enumerate}

\subsection{14. Устная версия ответа на 5
минут}\label{ux443ux441ux442ux43dux430ux44f-ux432ux435ux440ux441ux438ux44f-ux43eux442ux432ux435ux442ux430-ux43dux430-5-ux43cux438ux43dux443ux442}

Сначала формулирую абсолютную непрерывность мер:

\[
\nu\ll\mu
\quad\Longleftrightarrow\quad
\mu(A)=0\Rightarrow\nu(A)=0.
\]

Затем формулирую теорему Радона-Никодима: если \(\mu\)
\(\sigma\)-конечна, \(\nu\ll\mu\), то существует измеримая функция
\(f\), такая что

\[
\nu(A)=\int_A f\,d\mu.
\]

Эта функция называется производной Радона-Никодима:

\[
f=\frac{d\nu}{d\mu},
\]

и она единственна \(\mu\)-почти всюду.

Дальше объясняю плотности. Если \(X\) - случайный вектор в
\(\mathbb R^n\) и

\[
P_X\ll\lambda_n,
\]

то

\[
p_X=\frac{dP_X}{d\lambda_n}
\]

и

\[
P\{X\in B\}=\int_Bp_X(x)\,dx.
\]

Если абсолютной непрерывности нет, плотности относительно Лебега нет.
Например, у константы \(X=a\) распределение \(\delta_a\), и оно не имеет
лебеговой плотности.

Потом перехожу к условному ожиданию. Пусть \(X\in L^1\),
\(\mathcal G\subset\mathcal F\). Случайная величина \(Y\) называется
\(\mathbb E[X\mid\mathcal G]\), если \(Y\) \(\mathcal G\)-измерима и для
любого \(G\in\mathcal G\)

\[
\int_GY\,dP=\int_GX\,dP.
\]

Существование строится через Радона-Никодима: задаем на \(\mathcal G\)

\[
Q(G)=\int_GX\,dP.
\]

Тогда \(Q\ll P|_{\mathcal G}\), поэтому

\[
\mathbb E[X\mid\mathcal G]
=
\frac{dQ}{dP|_{\mathcal G}}.
\]

Условная вероятность - частный случай:

\[
P(A\mid\mathcal G)=\mathbb E[\mathbf 1_A\mid\mathcal G].
\]

Называю свойства:

\[
\mathbb E\,\mathbb E[X\mid\mathcal G]=\mathbb EX,
\]

если \(X\) \(\mathcal G\)-измерима, то

\[
\mathbb E[X\mid\mathcal G]=X,
\]

если \(X\) независима от \(\mathcal G\), то

\[
\mathbb E[X\mid\mathcal G]=\mathbb EX,
\]

и башня:

\[
\mathcal H\subset\mathcal G
\Rightarrow
\mathbb E[\mathbb E[X\mid\mathcal G]\mid\mathcal H]
=
\mathbb E[X\mid\mathcal H].
\]

В конце даю вычислительный пример: если есть совместная плотность
\(p_{X,Y}\), то

\[
p_{X\mid Y}(x\mid y)=\frac{p_{X,Y}(x,y)}{p_Y(y)}
\]

и

\[
\mathbb E[X\mid Y=y]=\int x p_{X\mid Y}(x\mid y)\,dx.
\]

Главные ловушки: плотность существует не всегда; условное ожидание -
случайная величина, а не число; все версии определены почти наверное.

\subsection{15. Если осталось 2 минуты перед
экзаменом}\label{ux435ux441ux43bux438-ux43eux441ux442ux430ux43bux43eux441ux44c-2-ux43cux438ux43dux443ux442ux44b-ux43fux435ux440ux435ux434-ux44dux43aux437ux430ux43cux435ux43dux43eux43c}

Запомнить четыре формулы.

\textbf{Радон-Никодим:}

\[
\nu\ll\mu
\Rightarrow
\nu(A)=\int_A\frac{d\nu}{d\mu}\,d\mu.
\]

\textbf{Плотность распределения:}

\[
P_X\ll\lambda_n,
\qquad
p_X=\frac{dP_X}{d\lambda_n},
\qquad
P_X(B)=\int_Bp_X\,dx.
\]

\textbf{Условное ожидание:}

\[
Y=\mathbb E[X\mid\mathcal G]
\iff
Y\text{ }\mathcal G\text{-измерима и }
\int_GY\,dP=\int_GX\,dP.
\]

\textbf{Построение через Радона-Никодима:}

\[
Q(G)=\int_GX\,dP,
\qquad
\mathbb E[X\mid\mathcal G]=\frac{dQ}{dP|_{\mathcal G}}.
\]

Главная ловушка: условное ожидание и плотность определяются почти всюду;
поточечные значения на нулевых множествах не важны.

\newpage

\end{document}
